\begin{lemma}
Let $p,k$ be natural numbers with $p\geq1$ and $k\geq1$.  The number of
ambient vector tuples
\[
 (\mathbf a_1,\mathbf b_1,\ldots,\mathbf a_k,\mathbf b_k)
 \in(\mathbb F_2^p)^{2k}
\]
such that
\[
 \langle\mathbf a_i,\mathbf b_j\rangle=1
 \qquad\text{for all }1\leq i\leq j\leq k
\]
is at most
\[
 \sum_{t=1}^{p}\binom kt\cdot
 2^{p(t+k)-\binom{t+1}{2}}.
\]
This is an intentional overcount in the application to projective
vectors: it includes zero vectors and does not quotient by the projective
equivalence relation.
\end{lemma}
\begin{proof}
Because $k\geq1$, the first diagonal equation is present.  For an ambient
tuple define $r_0=0$ and
$r_i=\dim\operatorname{span}(\mathbf a_1,\ldots,\mathbf a_i)$.  The first
diagonal equation forces $r_1=1$, and each subsequent rank is either the
previous rank or one larger.  If the final rank is $t$, the set of indices
at which the rank rises has at most $\binom kt$ choices.

Fix one such rank sequence.  At a rank-raising index there are at most
$2^p$ choices for $\mathbf a_i$, and the equations against the span give at
most $2^{p-r_i}$ choices for $\mathbf b_i$.  At a non-raising index there
are at most $2^{r_i}$ choices for $\mathbf a_i$, and again at most
$2^{p-r_i}$ choices for $\mathbf b_i$.  Reading the raising indices in
increasing order, their contribution is
\[
 \prod_{i=1}^{t}2^{2p-i},
\]
and every other index contributes $2^p$.  Thus this rank sequence gives at
most
\[
 \left(\prod_{i=1}^{t}2^{2p-i}\right)2^{p(k-t)}
 =2^{pt+pk-\binom{t+1}{2}}
\]
tuples.  Summing over the at most $\binom kt$ sequences and over
$1\leq t\leq p$ proves the formula.
\end{proof}
